\documentclass[aspectratio=169,10pt]{beamer}

% Shared Beamer setup for Fall 2026 course slides.
% Clean Cal Poly-inspired theme: no custom class files or uncommon packages.
\mode<presentation>{
  \usetheme{default}
  \usefonttheme{professionalfonts}
}

\definecolor{CalPolyGreen}{HTML}{154734}
\definecolor{CalPolyGold}{HTML}{BD8B13}
\definecolor{CalPolySage}{HTML}{7FA28F}
\definecolor{CalPolyMint}{HTML}{DDF1E7}
\definecolor{CalPolySand}{HTML}{A29F86}
\definecolor{CalPolyBeige}{HTML}{EFEEDE}
\definecolor{DeckInk}{HTML}{1F2933}
\definecolor{DeckMuted}{HTML}{5F6B7A}
\definecolor{DeckSoft}{HTML}{F7F7F5}

\setbeamersize{text margin left=0.55in,text margin right=0.55in}
\setbeamercolor{normal text}{fg=DeckInk,bg=white}
\setbeamercolor{structure}{fg=CalPolyGreen}
\setbeamercolor{title}{fg=white}
\setbeamercolor{frametitle}{fg=CalPolyGreen,bg=white}
\setbeamercolor{block title}{bg=CalPolySage,fg=white}
\setbeamercolor{block body}{bg=CalPolyMint,fg=black}
\setbeamercolor{alerted text}{fg=CalPolyGold}

\setbeamerfont{title}{size=\Huge,series=\bfseries}
\setbeamerfont{frametitle}{size=\Large,series=\bfseries}
\setbeamerfont{block title}{size=\normalsize,series=\bfseries}

\setbeamertemplate{navigation symbols}{}
\setbeamertemplate{itemize item}{\color{CalPolyGold}$\blacktriangleright$}
\setbeamertemplate{itemize subitem}{\color{CalPolyGreen}$\bullet$}
\setbeamertemplate{footline}{
  \hfill{\color{DeckMuted}\scriptsize\insertframenumber/\inserttotalframenumber}
  \hspace{0.18in}\vspace{0.12in}
}
\setbeamertemplate{frametitle}{
  \vspace{0.18in}
  \hspace{0.02in}\insertframetitle\par
  \vspace{0.08in}
  \noindent\makebox[\linewidth][l]{%
    {\color{CalPolyGold}\rule{0.55in}{2pt}}%
    {\color{CalPolyGreen}\rule{\dimexpr\linewidth-0.55in\relax}{2pt}}%
  }
}

\newcommand{\CourseNumber}{Course}
\newcommand{\CourseTitle}{Course Title}
\newcommand{\LectureNumber}{0}
\newcommand{\LectureTitle}{Lecture Title}
\newcommand{\LectureDate}{Date}
\newcommand{\CourseUnit}{Unit}

\newcommand{\coursenumber}[1]{\renewcommand{\CourseNumber}{#1}}
\newcommand{\coursetitle}[1]{\renewcommand{\CourseTitle}{#1}}
\newcommand{\lecturenumber}[1]{\renewcommand{\LectureNumber}{#1}}
\newcommand{\lecturetitle}[1]{\renewcommand{\LectureTitle}{#1}}
\newcommand{\lecturedate}[1]{\renewcommand{\LectureDate}{#1}}
\newcommand{\courseunit}[1]{\renewcommand{\CourseUnit}{#1}}

\author{Kirk Duran}
\institute{California Polytechnic State University, San Luis Obispo}
\date{\LectureDate}

\newcommand{\makedecktitle}{
  \begingroup
  \setbeamercolor{background canvas}{bg=CalPolyGreen}
  \begin{frame}[plain]
    \vfill
    \begin{center}
      {\usebeamerfont{title}\color{white}\LectureTitle\par}
      \vspace{0.18in}
      {\Large\color{white}Lecture \LectureNumber\par}
    \end{center}
    \vfill
    \noindent{\color{white}\rule{\linewidth}{1.2pt}}\par
    \vspace{0.08in}
    \noindent\colorbox{CalPolyGold}{\makebox[0.24\linewidth][c]{\color{white}\Large\CourseNumber}}
    \hspace{0.12in}{\color{white}\Large Kirk Duran}
  \end{frame}
  \endgroup
}

\newenvironment{keyidea}{\begin{block}{Key idea}}{\end{block}}
\newenvironment{activity}{\begin{block}{In class}}{\end{block}}
\newenvironment{cpdefinition}{\begin{block}{Definition}}{\end{block}}
\newenvironment{exampleblockcp}{
  \begingroup
  \setbeamercolor{block title}{bg=CalPolySand,fg=white}
  \setbeamercolor{block body}{bg=CalPolyBeige,fg=black}
  \begin{block}{Example}
}{
  \end{block}
  \endgroup
}

\newcommand{\imageplaceholder}[2][0.3in]{%
  \noindent\makebox[\linewidth][c]{%
    \fcolorbox{CalPolySage}{CalPolyBeige}{%
      \parbox[c][#1][c]{0.86\linewidth}{%
        \centering
        {\color{DeckMuted}\scriptsize Image placeholder: }%
        {\color{DeckInk}\scriptsize\ttfamily\detokenize{#2}}%
      }%
    }%
  }%
}

\newcommand{\sectionframe}[1]{
  \begingroup
  \setbeamercolor{background canvas}{bg=CalPolyGreen}
  \begin{frame}[plain]
    \vfill
    {\color{white}\LARGE\bfseries #1\par}
    \vspace{0.18in}
    {\color{CalPolyGold}\rule{0.22\paperwidth}{3pt}}
    {\color{white}\rule{0.62\paperwidth}{3pt}}
    \vfill
  \end{frame}
  \endgroup
}

\newcommand{\placeholdernote}{
  \begin{block}{Placeholder}
    Replace these bullets with the final lecture material, examples, and in-class activities.
  \end{block}
}


\pdfstringdefDisableCommands{\def\translate#1{#1}}

\graphicspath{{../assets/lecture-16/}}

% If an image file is not yet available, the slide displays a labeled
% placeholder using the intended filename. Place the matching file in
% ../assets/lecture-16/ to replace the placeholder automatically.
\newcommand{\lectureimage}[2][1.75in]{%
  \IfFileExists{../assets/lecture-16/#2}{%
    \includegraphics[width=\linewidth,height=#1,keepaspectratio]{#2}%
  }{%
    \fbox{%
      \parbox[c][#1][c]{0.94\linewidth}{%
        \centering
        \small\textbf{Image Placeholder}\\[0.08in]
        \scriptsize\texttt{\detokenize{#2}}
      }%
    }%
  }%
}

\setbeamersize{text margin left=0.42in,text margin right=0.42in}

\coursenumber{CS 1001}
\coursetitle{Introduction to Computing}
\lecturenumber{16}
\lecturetitle{The Function Design Process}
\lecturedate{Fall 2026}
\courseunit{2. Function Development and Verification}

\begin{document}

\makedecktitle

\begin{frame}{Today}
  \begin{itemize}
    \item Use a six-step process to move from a problem statement to a tested function.
    \item Identify the input and output data before writing implementation code.
    \item State the function's purpose precisely enough to determine expected behavior.
    \item Write a typed function header that exposes the intended interface.
    \item Derive tests from the specification before implementing the body.
    \item Implement, run tests, inspect failures, and iterate until the behavior matches the specification.
  \end{itemize}

  \begin{keyidea}
    The function body is not the starting point. A disciplined design process determines what the function must do before deciding how it will do it.
  \end{keyidea}
\end{frame}

\sectionframe{Part 1: Why Use a Function Design Process?}

\begin{frame}{Coding Immediately Often Mixes Several Problems Together}
  \begin{block}{Prompt}
    Write a function that determines whether a Fahrenheit temperature is freezing.
  \end{block}

  \begin{columns}[T,onlytextwidth]
    \begin{column}{0.46\textwidth}
      \begin{block}{Immediate coding}
        \texttt{def freezing(temp):}\\
        \hspace*{0.20in}\texttt{...}
      \end{block}
    \end{column}
    \begin{column}{0.46\textwidth}
      \begin{block}{Questions still unanswered}
        What input type is expected?\\
        What value should be returned?\\
        Is exactly \texttt{32} freezing?\\
        How will we verify the boundary?
      \end{block}
    \end{column}
  \end{columns}

  \begin{keyidea}
    Implementation becomes easier when the data and expected behavior have already been made explicit.
  \end{keyidea}
\end{frame}

\begin{frame}{The Six-Step Function Design Process}
  \begin{enumerate}
    \item \textbf{Identify the data and types.}
    \item \textbf{Write a one-line purpose statement.}
    \item \textbf{Write the function header with type hints.}
    \item \textbf{Write tests from the expected behavior.}
    \item \textbf{Write the function body.}
    \item \textbf{Run, inspect, and iterate.}
  \end{enumerate}

  \begin{center}
    % Intended visual: six-stage pipeline from problem statement to verified function.
    \lectureimage[1.65in]{function-design-process.png}
  \end{center}
\end{frame}

\begin{frame}{Specification Comes Before Implementation}
  \begin{cpdefinition}
    A \textbf{specification} describes the behavior a function is expected to provide: what data it accepts and what result it should produce.
  \end{cpdefinition}

  \begin{block}{Two different questions}
    \textbf{Specification:} What should this function do?\\[0.08in]
    \textbf{Implementation:} How will the function produce that behavior?
  \end{block}

  \begin{keyidea}
    Tests should be derived from the specification. Otherwise a test may merely confirm what the current implementation already happens to do.
  \end{keyidea}
\end{frame}

\begin{frame}{Running Example: Determine Whether a Temperature Is Freezing}
  \begin{block}{Problem statement}
    Design a function that accepts a temperature in degrees Fahrenheit and returns whether that temperature is at or below the freezing point of water.
  \end{block}

  \begin{block}{Do not implement yet}
    We will deliberately postpone the function body until the input, output, purpose, interface, and tests are established.
  \end{block}

  \begin{activity}
    Before continuing, identify the input data, its likely Python type, and the type of the result.
  \end{activity}
\end{frame}

\sectionframe{Part 2: Step 1 --- Identify Data, Variables, and Types}

\begin{frame}{Begin with the Data Crossing the Function Boundary}
  \begin{columns}[T,onlytextwidth]
    \begin{column}{0.46\textwidth}
      \begin{block}{Input}
        Fahrenheit temperature\\[0.08in]
        suggested name: \texttt{temperature}\\[0.08in]
        expected type: \texttt{float}
      \end{block}
    \end{column}
    \begin{column}{0.46\textwidth}
      \begin{block}{Output}
        whether the temperature is freezing\\[0.08in]
        result type: \texttt{bool}
      \end{block}
    \end{column}
  \end{columns}

  \begin{keyidea}
    The input and output types describe the shape of the computation before any algorithm has been chosen.
  \end{keyidea}
\end{frame}

\begin{frame}{Choose Names That Describe the Role of the Data}
  \begin{columns}[T,onlytextwidth]
    \begin{column}{0.46\textwidth}
      \begin{block}{Weak names}
        \texttt{x}\\
        \texttt{num}\\
        \texttt{value}
      \end{block}
    \end{column}
    \begin{column}{0.46\textwidth}
      \begin{block}{More informative names}
        \texttt{temperature}\\
        \texttt{weight}\\
        \texttt{tax\_rate}
      \end{block}
    \end{column}
  \end{columns}

  \begin{itemize}
    \item A parameter name should communicate the role of the argument supplied by the caller.
    \item Good names reduce the amount of interpretation required when reading the function header and tests.
  \end{itemize}
\end{frame}

\begin{frame}{Type Hints Document Intended Data}
  \begin{block}{Example annotations}
    \texttt{temperature: float}\\
    \texttt{count: int}\\
    \texttt{name: str}\\
    \texttt{eligible: bool}
  \end{block}

  \begin{alertblock}{Technical distinction}
    Type hints are annotations. Ordinary Python execution does not generally reject a call merely because an argument disagrees with an annotation.
  \end{alertblock}

  \begin{keyidea}
    In this course, type hints serve as explicit documentation of the interface and support reasoning about valid inputs and expected outputs.
  \end{keyidea}
\end{frame}

\begin{frame}{Multiple Inputs Require Separate Roles}
  \begin{block}{Problem}
    Return the larger of two integers.
  \end{block}

  \begin{center}
    \renewcommand{\arraystretch}{1.35}
    \begin{tabular}{l|l|l}
      \textbf{Role} & \textbf{Name} & \textbf{Type} \\\hline
      first candidate & \texttt{a} & \texttt{int} \\
      second candidate & \texttt{b} & \texttt{int} \\
      result & --- & \texttt{int}
    \end{tabular}
  \end{center}

  \begin{keyidea}
    Step 1 is complete when the data entering and leaving the function is sufficiently explicit to describe the interface.
  \end{keyidea}
\end{frame}

\sectionframe{Part 3: Steps 2 and 3 --- Purpose and Function Header}

\begin{frame}{Step 2: Write a One-Line Purpose Statement}
  \begin{cpdefinition}
    A \textbf{purpose statement} is a concise description of the behavior the function provides to its caller.
  \end{cpdefinition}

  \begin{block}{Running example}
    Return \texttt{True} when \texttt{temperature} is at or below \texttt{32.0} degrees Fahrenheit; otherwise return \texttt{False}.
  \end{block}

  \begin{keyidea}
    A useful purpose statement is specific enough that another person can derive expected outputs without reading the implementation.
  \end{keyidea}
\end{frame}

\begin{frame}{Purpose Statements Describe Behavior, Not Syntax}
  \begin{columns}[T,onlytextwidth]
    \begin{column}{0.46\textwidth}
      \begin{alertblock}{Too vague}
        ``Checks the temperature.''
      \end{alertblock}
      \begin{alertblock}{Too implementation-specific}
        ``Uses \texttt{<=} to compare the argument with 32.''
      \end{alertblock}
    \end{column}
    \begin{column}{0.46\textwidth}
      \begin{exampleblockcp}
        ``Return whether a Fahrenheit temperature is at or below the freezing point of water.''
      \end{exampleblockcp}
    \end{column}
  \end{columns}

  \begin{keyidea}
    State observable behavior. Leave implementation choices for Step 5.
  \end{keyidea}
\end{frame}

\begin{frame}{Step 3: Derive the Function Header from the Data}
  \begin{block}{Data from Step 1}
    input: \texttt{temperature: float}\\
    output: \texttt{bool}
  \end{block}

  \begin{block}{Function header}
    \centering
    \Large\texttt{def is\_freezing(temperature: float) -> bool:}
  \end{block}

  \begin{itemize}
    \item The function name summarizes the operation.
    \item The parameter names the incoming data.
    \item The parameter annotation records the intended input type.
    \item The return annotation records the intended result type.
  \end{itemize}
\end{frame}

\begin{frame}{The Header Exposes the Function's Interface}
  \begin{cpdefinition}
    A function's \textbf{interface} is the information a caller needs in order to use the function correctly, without depending on its internal implementation.
  \end{cpdefinition}

  \begin{center}
    % Intended visual: caller on left, typed function header in center, result on right; body hidden.
    \lectureimage[1.15in]{function-interface-boundary.png}
  \end{center}

  \begin{keyidea}
    The caller should be able to reason from the function's name, parameters, documented behavior, and returned result without tracing its body. Returning a different kind of result would violate that intended interface.
  \end{keyidea}
\end{frame}

\begin{frame}{Checkpoint: Can the Behavior Be Predicted from the Design?}
  \begin{block}{Current design}
    \texttt{def is\_freezing(temperature: float) -> bool:}\\[0.08in]
    Purpose: return whether the Fahrenheit temperature is at or below \texttt{32.0}.
  \end{block}

  \begin{activity}
    Without writing a body, determine the expected results for \texttt{31.0}, \texttt{32.0}, and \texttt{33.0}. Explain which value is the boundary case.
  \end{activity}

  \begin{keyidea}
    If expected results cannot be determined, the specification is not yet precise enough to support reliable testing.
  \end{keyidea}
\end{frame}

\sectionframe{Part 4: Step 4 --- Write Tests Before the Body}

\begin{frame}{Tests Translate the Specification into Executable Expectations}
  \begin{block}{Specification}
    Return whether \texttt{temperature} is at or below \texttt{32.0}.
  \end{block}

  \begin{center}
    \renewcommand{\arraystretch}{1.4}
    \begin{tabular}{c|c|l}
      \textbf{Input} & \textbf{Expected result} & \textbf{Reason} \\\hline
      \texttt{31.0} & \texttt{True} & below boundary \\
      \texttt{32.0} & \texttt{True} & exactly at boundary \\
      \texttt{33.0} & \texttt{False} & above boundary
    \end{tabular}
  \end{center}

  \begin{keyidea}
    The expected values come from the purpose statement, not from executing the unfinished function.
  \end{keyidea}
\end{frame}

\begin{frame}{Boundary Cases Are Determined by the Specification}
  \begin{block}{Decision boundary}
    \centering
    \Large\texttt{temperature <= 32.0}
  \end{block}

  \begin{itemize}
    \item A value below the boundary exercises one side of the decision.
    \item The boundary value checks whether equality is included.
    \item A value above the boundary exercises the opposite side.
  \end{itemize}

  \begin{center}
    % Intended visual: number line highlighting 31, 32, and 33 around the boundary.
    \lectureimage[1.65in]{boundary-test-cases.png}
  \end{center}
\end{frame}

\begin{frame}{Write the Test File Before Implementing the Body}
  \begin{block}{\texttt{test\_temperature.py}}
    \texttt{import unittest}\\
    \texttt{from temperature import is\_freezing}\\[0.08in]
    \texttt{class TestFreezing(unittest.TestCase):}\\
    \hspace*{0.20in}\texttt{def test\_below(self):}\\
    \hspace*{0.40in}\texttt{self.assertTrue(is\_freezing(31.0))}\\[0.05in]
    \hspace*{0.20in}\texttt{def test\_boundary(self):}\\
    \hspace*{0.40in}\texttt{self.assertTrue(is\_freezing(32.0))}\\[0.05in]
    \hspace*{0.20in}\texttt{def test\_above(self):}\\
    \hspace*{0.40in}\texttt{self.assertFalse(is\_freezing(33.0))}
  \end{block}

  \begin{keyidea}
    Before the body is complete, failing tests are expected: they establish the behavior the implementation must eventually satisfy.
  \end{keyidea}
\end{frame}

\begin{frame}{Separate Files Separate Responsibilities}
  \begin{columns}[T,onlytextwidth]
    \begin{column}{0.44\textwidth}
      \begin{block}{\texttt{temperature.py}}
        Defines the behavior being implemented.
      \end{block}
    \end{column}
    \begin{column}{0.44\textwidth}
      \begin{block}{\texttt{test\_temperature.py}}
        Defines examples of expected behavior.
      \end{block}
    \end{column}
  \end{columns}

  \begin{center}
    \texttt{from temperature import is\_freezing}
  \end{center}

  \begin{keyidea}
    The implementation and its tests evolve together, but they serve different purposes and should remain conceptually distinct.
  \end{keyidea}
\end{frame}

\begin{frame}{A Test Should Explain What Behavior It Is Checking}
  \begin{columns}[T,onlytextwidth]
    \begin{column}{0.46\textwidth}
      \begin{alertblock}{Weak}
        \texttt{def test1(self):}\\
        \hspace*{0.20in}\texttt{...}
      \end{alertblock}
    \end{column}
    \begin{column}{0.46\textwidth}
      \begin{exampleblockcp}
        \texttt{def test\_freezing\_point(self):}\\
        \hspace*{0.20in}\texttt{self.assertTrue(}\\
        \hspace*{0.40in}\texttt{is\_freezing(32.0))}
      \end{exampleblockcp}
    \end{column}
  \end{columns}

  \begin{itemize}
    \item The test name should identify the behavior or case under examination.
    \item A focused failure message is more informative than a large test containing unrelated expectations.
  \end{itemize}
\end{frame}

\begin{frame}{Tests Must Be Independent of the Implementation Strategy}
  \begin{block}{Acceptable expectation}
    \texttt{self.assertTrue(is\_freezing(32.0))}
  \end{block}

  \begin{alertblock}{Poor testing idea}
    ``Verify that the function uses an \texttt{if} statement.''
  \end{alertblock}

  \begin{itemize}
    \item The specification requires a Boolean result.
    \item It does not require one particular arrangement of statements to obtain that result.
  \end{itemize}

  \begin{keyidea}
    Test externally observable behavior unless the specification explicitly requires an implementation property.
  \end{keyidea}
\end{frame}

\sectionframe{Part 5: Step 5 --- Write the Function Body}

\begin{frame}{Only Now Do We Implement the Computation}
  \begin{block}{Design already established}
    \texttt{def is\_freezing(temperature: float) -> bool:}
  \end{block}

  \begin{block}{Implementation}
    \texttt{def is\_freezing(temperature: float) -> bool:}\\
    \hspace*{0.20in}\texttt{return temperature <= 32.0}
  \end{block}

  \begin{keyidea}
    The implementation is short because most of the reasoning has already been expressed in the specification, interface, and tests.
  \end{keyidea}
\end{frame}

\begin{frame}{Prefer the Simplest Body That Clearly Satisfies the Specification}
  \begin{columns}[T,onlytextwidth]
    \begin{column}{0.46\textwidth}
      \begin{block}{Direct predicate}
        \texttt{def is\_freezing(t: float) -> bool:}\\
        \hspace*{0.20in}\texttt{return t <= 32.0}
      \end{block}
    \end{column}
    \begin{column}{0.46\textwidth}
      \begin{block}{Equivalent conditional}
        \texttt{def is\_freezing(t: float) -> bool:}\\
        \hspace*{0.20in}\texttt{if t <= 32.0:}\\
        \hspace*{0.40in}\texttt{return True}\\
        \hspace*{0.20in}\texttt{return False}
      \end{block}
    \end{column}
  \end{columns}

  \begin{keyidea}
    Both implementations satisfy the same interface. Prefer the clearer expression of the underlying computation.
  \end{keyidea}
\end{frame}

\begin{frame}{Some Specifications Naturally Require Conditional Structure}
  \begin{block}{Problem}
    Return a shipping cost of \texttt{5.0} for packages weighing at most \texttt{10.0} pounds; otherwise return \texttt{8.0}.
  \end{block}

  \begin{block}{Implementation}
    \texttt{def shipping\_cost(weight: float) -> float:}\\
    \hspace*{0.20in}\texttt{if weight <= 10.0:}\\
    \hspace*{0.40in}\texttt{return 5.0}\\
    \hspace*{0.20in}\texttt{return 8.0}
  \end{block}

  \begin{keyidea}
    The design process does not prescribe one body shape; the specification determines what behavior must be represented.
  \end{keyidea}
\end{frame}

\begin{frame}{Trace the Body Against a Previously Chosen Test}
  \begin{block}{Test case}
    \texttt{shipping\_cost(10.0)} should return \texttt{5.0}.
  \end{block}

  \begin{enumerate}
    \item Bind \texttt{weight} to \texttt{10.0}.
    \item Evaluate \texttt{weight <= 10.0} $\rightarrow$ \texttt{True}.
    \item Enter the selected branch.
    \item Execute \texttt{return 5.0}.
    \item Compare actual result \texttt{5.0} with expected result \texttt{5.0}.
  \end{enumerate}

  \begin{keyidea}
    Tracing remains useful: tests tell us whether behavior matches; tracing explains how a particular result was produced.
  \end{keyidea}
\end{frame}

\begin{frame}{Do Not Implement Only the Examples}
  \begin{alertblock}{Incorrect implementation}
    \texttt{def is\_freezing(t: float) -> bool:}\\
    \hspace*{0.20in}\texttt{return t == 31.0 or t == 32.0}
  \end{alertblock}

  \begin{itemize}
    \item This passes two known positive examples.
    \item It does not implement the general rule in the specification.
    \item For example, \texttt{is\_freezing(0.0)} would be incorrect.
  \end{itemize}

  \begin{keyidea}
    Tests are examples of the specification; they are not a replacement for understanding the general behavior being specified.
  \end{keyidea}
\end{frame}

\sectionframe{Part 6: Step 6 --- Run, Inspect, and Iterate}

\begin{frame}{Run the Test Suite After Implementing the Body}
  \begin{center}
    % Intended visual: test suite feeding implementation and producing pass/fail result.
    \lectureimage[1.55in]{run-test-suite.png}
  \end{center}

  \begin{columns}[T,onlytextwidth]
    \begin{column}{0.44\textwidth}
      \begin{block}{Green}
        Current implementation satisfies the tested expectations.
      \end{block}
    \end{column}
    \begin{column}{0.44\textwidth}
      \begin{block}{Red}
        At least one expectation is not satisfied or a test cannot complete normally.
      \end{block}
    \end{column}
  \end{columns}
\end{frame}

\begin{frame}{A Failure Creates a New Reasoning Problem}
  \begin{block}{Suppose the body is}
    \texttt{def is\_freezing(t: float) -> bool:}\\
    \hspace*{0.20in}\texttt{return t < 32.0}
  \end{block}

  \begin{block}{Observed result}
    \texttt{test\_below} passes.\\
    \texttt{test\_boundary} fails.\\
    \texttt{test\_above} passes.
  \end{block}

  \begin{activity}
    What does this failure pattern tell you about the relationship between the implementation and the specification?
  \end{activity}
\end{frame}

\begin{frame}{Debug the First Meaningful Divergence}
  \begin{enumerate}
    \item Start from the failing boundary test.
    \item Expected condition result for \texttt{32.0}: \texttt{True}.
    \item Actual condition result for \texttt{32.0 < 32.0}: \texttt{False}.
    \item The divergence occurs in the comparison operator.
    \item Correct \texttt{<} to \texttt{<=}.
    \item Run the complete test suite again.
  \end{enumerate}

  \begin{keyidea}
    Lecture 15's debugging process is now embedded inside the larger function design process.
  \end{keyidea}
\end{frame}

\begin{frame}{Iteration Can Improve Either Code or Specification}
  \begin{columns}[T,onlytextwidth]
    \begin{column}{0.45\textwidth}
      \begin{block}{Implementation problem}
        The stated behavior is clear, but the body does not produce it.
      \end{block}
    \end{column}
    \begin{column}{0.45\textwidth}
      \begin{block}{Specification problem}
        A reasonable input exposes behavior that the original statement never defined.
      \end{block}
    \end{column}
  \end{columns}

  \begin{block}{Example question}
    What should \texttt{shipping\_cost(-2.0)} mean? The original specification may need to state whether negative weights are valid inputs.
  \end{block}
\end{frame}

\begin{frame}{Passing Tests Should Remain as Regression Tests}
  \begin{cpdefinition}
    A \textbf{regression} occurs when a change causes behavior that previously worked to become incorrect.
  \end{cpdefinition}

  \begin{itemize}
    \item Keep existing tests when the implementation changes.
    \item Add a new test when a newly discovered case should remain supported.
    \item Run the whole test suite, not only the test related to the latest edit.
  \end{itemize}

  \begin{keyidea}
    A growing test suite preserves executable examples of behavior that future changes should continue to satisfy.
  \end{keyidea}
\end{frame}

\begin{frame}{The Complete Development Loop}
  \begin{center}
    % Intended visual: specification -> header -> tests -> red -> body -> green -> inspect/refine -> tests.
    \lectureimage[2.75in]{design-test-iterate-loop.png}
  \end{center}
\end{frame}

\begin{frame}[shrink=5]{Worked Example: Return the Larger of Two Integers}
  \begin{enumerate}
    \item \textbf{Data:} \texttt{a: int}, \texttt{b: int}; result \texttt{int}.
    \item \textbf{Purpose:} return the larger of \texttt{a} and \texttt{b}; if they are equal, return that value.
    \item \textbf{Header:} \texttt{def maximum(a: int, b: int) -> int:}
    \item \textbf{Tests:}
      \begin{itemize}
        \item \texttt{maximum(5, 3)} $\rightarrow$ \texttt{5}
        \item \texttt{maximum(3, 5)} $\rightarrow$ \texttt{5}
        \item \texttt{maximum(5, 5)} $\rightarrow$ \texttt{5}
      \end{itemize}
    \item \textbf{Body:}
      \texttt{if a >= b: return a}\quad otherwise \texttt{return b}
    \item \textbf{Iterate:} run all three tests and inspect any mismatch.
  \end{enumerate}

  \begin{keyidea}
    The equality case is deliberate: it tests the boundary between ``first is larger'' and ``second is larger.''
  \end{keyidea}
\end{frame}

\begin{frame}[shrink=4]{Integrated Exercise: Design Before Coding}
  \begin{block}{Problem}
    Design \texttt{has\_minimum\_length(password: str) -> bool}, which returns whether \texttt{password} contains at least eight characters.
  \end{block}

  \begin{activity}
    Before writing the body, produce:
    \begin{enumerate}
      \item the input and output data;
      \item a one-line purpose statement;
      \item the complete typed header;
      \item three tests that examine the boundary;
      \item only then, the simplest correct body.
    \end{enumerate}
  \end{activity}
\end{frame}

\begin{frame}{A Practical Checklist for New Functions}
  \begin{enumerate}
    \item What values enter the function, and what are their intended types?
    \item What value must the function return?
    \item Can the behavior be stated precisely in one sentence?
    \item Does the function header reflect that data and behavior?
    \item Which ordinary, boundary, and special cases follow from the specification?
    \item Do the tests encode those expectations before the body is trusted?
    \item Does the simplest clear implementation satisfy the tests?
    \item After a correction, did the entire suite still pass?
  \end{enumerate}
\end{frame}

\begin{frame}{Takeaways}
  \begin{itemize}
    \item Function design begins with data and behavior, not implementation syntax.
    \item Type hints document the intended interface but are not, by themselves, runtime enforcement.
    \item Purpose statements make expected behavior explicit enough to test.
    \item Tests should be derived from the specification before implementation details influence expectations.
    \item The body implements the general behavior; it should not merely memorize the examples in the test suite.
    \item Red--green feedback and debugging turn failures into evidence for the next iteration.
  \end{itemize}

  \begin{keyidea}
    Design, testing, implementation, and debugging are not separate habits. Together they form a repeatable process for producing functions whose behavior can be explained and verified.
  \end{keyidea}
\end{frame}

\end{document}
